\section{The Proposed Framework: MAP}
\label{sec:method}

In this section, we introduce our framework, MAP.
The framework of MAP is shown in Fig. \ref{fig:framework}.
We first revisit the notion of matrix adaptation from a vector space perspective.
As discussed in the introduction, any weight matrix $\mathbf{W} \in \mathbb{R}^{n \times m}$ can be flattened into a high-dimensional vector $\mathbf{w} \in \mathbb{R}^{nm}$. 
This allows us to interpret the adaptation process in terms of classical vector operations, where the direction of $\mathbf{w}$ is defined by its unit-norm vector, and the magnitude by its $\ell_2$ norm.


\subsection{Vector-Based Formulation}

Under this formulation, we model the final fine-tuned weight vector $\mathbf{w}'$ as a directional combination of the pre-trained weights and the learned update vector:
\begin{equation}
\mathbf{w}' = \alpha \cdot \frac{\mathbf{w}}{\|\mathbf{w}\|_2+\epsilon} + \beta \cdot \frac{\Delta\mathbf{w}}{\|\Delta\mathbf{w}\|_2+\epsilon},
\label{eq:vector-map}
\end{equation}
where $\alpha, \beta \in \mathbb{R}$ are learnable scalar coefficients that independently control the contribution (i.e., magnitude) of the pre-trained vector and its directional update, and $\epsilon>0$ is a small constant for numerical stability.
This formulation has several appealing properties:
\begin{itemize}
    \item It cleanly separates the direction and magnitude of each component.
    \item It treats $\mathbf{w}$ and $\Delta \mathbf{w}$ in the same vector space, enabling their alignment or contrast to be interpreted geometrically.
    \item It introduces only two additional parameters per layer, making it highly lightweight compared to methods like DoRA.


\end{itemize}


\subsection{Matrix-Based Formulation}

When applied to the original matrix space, particularly when integrated with LoRA, MAP takes the following form:
\begin{equation}
\mathbf{W}^* = \alpha \cdot \frac{\mathbf{W}}{\|\mathbf{W}\|_F+\epsilon} + \beta \cdot \frac{\mathbf{A}\mathbf{B}}{\|\mathbf{A}\mathbf{B}\|_F+\epsilon},
\label{eq:matrix-map}
\end{equation}
where $\|\cdot\|_F$ denotes the Frobenius norm, and $\Delta \mathbf{W} = \mathbf{A}\mathbf{B}$ is the standard low-rank adaptation used in LoRA.
This formulation can be directly derived since the Frobenius norm of $\mathbf{W}$ satisfies:
\begin{equation}
\|\mathbf{W}\|_F = \|\mathbf{w}\|_2,
\end{equation}
Therefore, the vector and matrix formulations of MAP are mathematically equivalent in terms of magnitude scaling.

\subsection{Numerical Stability and Initialization}

A direct implementation of Eq. \ref{eq:matrix-map} must handle the standard LoRA initialization in which $\mathbf{A}$ is initialized with a Kaiming distribution and $\mathbf{B}$ is initialized to zero.
At this point $\mathbf{A}\mathbf{B}=\mathbf{0}$, making normalization by $\|\mathbf{A}\mathbf{B}\|_F$ undefined.
We therefore use the $\epsilon$-stabilized normalization in Eq. \ref{eq:matrix-map}, with $\epsilon=10^{-6}$ in our implementation.
Unless otherwise specified, the scalar coefficients are initialized as
\begin{equation}
\alpha_0=\|\mathbf{W}\|_F,\qquad \beta_0=1.
\end{equation}
This keeps the initial base branch close to the pre-trained weight,
$\alpha_0\mathbf{W}/(\|\mathbf{W}\|_F+\epsilon)\approx \mathbf{W}$,
while the update branch is initially zero-valued because $\mathbf{A}\mathbf{B}=\mathbf{0}$.
The denominator is optimized jointly with the low-rank matrices rather than detached.
This design preserves the standard LoRA initialization while making MAP well-defined from the first optimization step.

\paragraph{Initialization choices.}
The default of $\alpha_0=\|\mathbf{W}\|_F$ and $\beta_0=1$ is \emph{function-preserving} up to $O(\epsilon)$: the network output is identical to the frozen pre-trained model at $t=0$, while gradients with respect to $\alpha$, $\beta$, $\mathbf{A}$, and $\mathbf{B}$ are all well-defined.
Alternative choices include $\beta_0=0$ (fully function-preserving but suppresses early update-direction learning) and $\beta_0\ll 1$ (a warm-up compromise).
The denominator $\|\mathbf{A}\mathbf{B}\|_F+\epsilon$ is optimized jointly with $\mathbf{A}$ and $\mathbf{B}$ rather than detached; detaching the denominator decouples magnitude and direction optimization more aggressively but can slow convergence.
We report sensitivity analyses for $\epsilon$, $\beta_0$, and the detach choice in the supplementary material.

\subsection{Geometric Properties}

The Frobenius formulation gives MAP invariance properties that are not shared by column-wise normalizations.
For any orthogonal matrices $\mathbf{U}$ and $\mathbf{V}$,
\begin{equation}
\|\mathbf{U}\mathbf{W}\mathbf{V}\|_F^2
=\mathrm{tr}(\mathbf{V}^{\top}\mathbf{W}^{\top}\mathbf{U}^{\top}\mathbf{U}\mathbf{W}\mathbf{V})
=\|\mathbf{W}\|_F^2.
\end{equation}
Thus the global direction $\mathbf{W}/\|\mathbf{W}\|_F$ is preserved under left and right orthogonal changes of basis, whereas column-wise normalization is not invariant to right-side column mixing.

The normalization also clarifies the optimization geometry.
Let $\Delta=\mathbf{A}\mathbf{B}$ and, for nonzero $\Delta$, define $\hat{\delta}=\mathrm{vec}(\Delta)/\|\Delta\|_F$.
The first-order differential of the normalized update direction is
\begin{equation}
d\hat{\delta}=
\frac{1}{\|\Delta\|_F}
\left(\mathbf{I}-\hat{\delta}\hat{\delta}^{\top}\right)d\,\mathrm{vec}(\Delta),
\end{equation}
which projects changes in $\Delta$ onto the tangent space of the Frobenius sphere.
Consequently, $\beta$ controls update magnitude, while $\mathbf{A}$ and $\mathbf{B}$ primarily shape the update direction.

MAP can also be written as an effective update over the pre-trained weight:
\begin{equation}
\mathbf{W}^*-\mathbf{W}
=
\left(\frac{\alpha}{\|\mathbf{W}\|_F+\epsilon}-1\right)\mathbf{W}
+
\beta\frac{\mathbf{A}\mathbf{B}}{\|\mathbf{A}\mathbf{B}\|_F+\epsilon}.
\end{equation}
This decomposition shows that MAP augments the low-rank update with a full-rank radial component along the pre-trained weight direction.
Ignoring the negligible $\epsilon$ term, the distance to the pre-trained weight is bounded by
\begin{equation}
\|\mathbf{W}^*-\mathbf{W}\|_F
\leq
|\alpha-\|\mathbf{W}\|_F|+|\beta|,
\end{equation}
so the learned scalars provide direct control over the adaptation radius.

\input{sec/tables/cr}


In this way, MAP can be viewed as a direction-aware modulation layer that scales the normalized base and update matrices in a principled and learnable fashion.
We advocate LoMAP as a drop-in enhancement to existing PEFT frameworks.
In our experiments, we integrate MAP with LoRA to form a new variant, LoMAP, which we adopt as the default implementation.
